\section*{Acknowledgments}
We would like to thank X. Ji,  Y. Ji, C. Monahan, T.  Neumann, A. Shindler, V. Harlander, F. Lange and P. Petreczky for useful discussions. We thank A. Vairo for reading the paper and giving very useful comments. The work of N.~B. and X.-P.~W. is supported by the DFG (Deutsche Forschungsgemeinschaft,
German Research Foundation) Grant No. BR 4058/2-2. We
acknowledge support from the DFG cluster of excellence ``ORIGINS'' under
Germany's Excellence Strategy - EXC-2094 - 390783311. The authors acknowledge support from STRONG-2020- European Union’s Horizon 2020 research and innovation program under grant agreement
No. 824093.

\appendix
\section{Quark quasi-PDF operator  in gradient flow}
\label{app:quarkquasi}
\begin{figure}
\centering
\includegraphics[width =0.72\textwidth]{images/quarkquasi.pdf}
\caption{Feynman diagrams  that contribute to the one-loop corrections for the hadronic matrix element defined as eq.~(\ref{eq:quarkhmte}) in gradient-flow formalism. The dashed line is the Wilson line, the filled squares represent the flowed $B_{\mu}$ field at flow time $t$, the cross circles represent the spacetime points that connecting the Wilson line and the fermion fields, which are all flowed at the flow time $t$. Diagrams that give contribution proportional to the external momentum and thus vanish in our calculation are dropped. The external quark one-loop self-energy diagrams are not listed here, they are shown in figure.~\ref{fig:quarkself}}
\label{fig:sqpdf1}
\end{figure}
In this appendix, we consider the hadronic matrix element constructed using the operator defined in eq.~(\ref{eq:defqoperator}),
\begin{eqnarray}
\label{eq:quarkhmte}
h(z, P^z)= \langle P| \bar{\psi}(zv)\Gamma W(zv,0)\psi(0)| P\rangle,
\end{eqnarray}
where $P$ is the hadronic state with momentum $P^z$, $\Gamma = \gamma\cdot v$ or $\Gamma= \gamma^{\alpha}_{\perp}$.
The one-loop matching calculation for $h(z, P^z)$ from the gradient-flow scheme to the $\overline{\rm MS}$-like schemes was done in ref.~\cite{Monahan:2017hpu} with $P^z=0$ and full flow time dependence. 

However, the matching coefficient given by eq.~(44) of ref.~\cite{Monahan:2017hpu} does not lead to $z$ independent result in the small flow-time limit, 
which contradicts with our argument that we only need to do matching calculations for the local current operators stemming from the Wilson-line operators based on the one-dimensional auxiliary-field formalism.
Therefore, we re-calculate the matching coefficient in great detail and compare with the calculation done in ref.~\cite{Monahan:2017hpu}\footnote{Our results correct and supersede the results of ref.~\cite{Monahan:2017hpu}, for which an erratum is in preparation.}. 

The matching relation for $h(z, P^z)$ can be expressed as
\begin{eqnarray}
h^{\rm R}(z, t) = \mathcal{C}_q(t, z, \mu)h^{\overline{\rm MS}}(z, \mu),
\end{eqnarray}
 where we have dropped the exponentiated ``mass” correction $\delta m$ in the matching relation and choose $P^z =0$ for simplicity and comparison. 
 
 The Feynman diagrams that contribute to the  one-loop corrections for $h(z, P^z)$ in gradient-flow formalism are shown in figure.~\ref{fig:quarkself} and  figure.~\ref{fig:sqpdf1}.
The contribution coming from the quark self-energy diagrams in figure.~\ref{fig:quarkself} can be summarized as  
\begin{eqnarray}
\label{eq:appquarkself}
\mathcal{M}_{\mathring{\chi}}^{\rm R} =  \frac{\alpha_s}{4\pi}C_F\left[\frac{1}{\epsilon_{\rm IR}}+\log\left(2t\mu^2\right) + \gamma_E - \log(432) \right],
\end{eqnarray}
where the subscript $\mathring{\chi}$ indicates that we have used the ringed fermion field to remove the UV divergences.
%%%%%%%%%%%%%%%%%%%%%%%%%%%%%%%%%%%%%%%%%%%%%%%%%
\subsection{Wilson line self-energy diagram $(f_1)$}
 The one-loop Wilson line self energy diagram is shown as diagram $(f_1)$ in figure.~\ref{fig:sqpdf1} and its amplitude is given by
 \begin{eqnarray}
\mathcal{M}_{(f_1)}&=&-g^2\tilde{\mu}^{2\epsilon} C_F \int_0^z d\, s_1\int_0^{s_1} d\, s_2\, \int \frac{d^d k}{(2\pi)^d} \frac{1}{k^2}e^{-2t k^2} e^{i(k\cdot v)(s_2-s_1)},
\end{eqnarray}
which is finite at $d=4$. Therefore, we set $d=4$ and obtain,
\begin{eqnarray}
\label{eq:Wilsonselffull}
\mathcal{M}_{(f_1)}&=&-g^2C_F \int_0^z d\, s_1\int_0^{s_1} d\, s_2\,\int_0^{+\infty} dx \frac{1}{(2\pi)^4}e^{-\frac{(s_2-s_1)^2}{4(2t +x)}}\left(\frac{\pi}{2t +x}\right)^{2}\nonumber\\
&=&\frac{\alpha_s}{4\pi}C_F\times 2\left[\log\left(\bar{z}^2\right) +\gamma_E+2\left(1-e^{-\bar{z}^2}\right)-\text{Ei}\left(-\bar{z}^2\right)-2\sqrt{\pi}\bar{z}\, \text{erf}\left(\bar{z}\right)\right],
\end{eqnarray}
where $\bar{z} =  \frac{z}{r_F}$ with $r_F =\sqrt{8t}$ representing the flow radius, $\text{Ei}\left(x\right)$ is the exponential integral
\begin{eqnarray}
\text{Ei}\left(x\right) = -\int_{-x}^{+\infty} \frac{e^{-t}}{t}dt,
\end{eqnarray}
and $\text{erf}\left(\bar{z}\right)$ is the error function
\begin{eqnarray}
\text{erf}\left(\bar{z}\right) = \frac{2}{\sqrt{\pi}}\int_0^{\bar{z}} e^{-t^2} dt.
\end{eqnarray}
In the small flow-time limit, we have
\begin{eqnarray}
\label{eq:selfsmallflowlimit}
\mathcal{M}_{(f_1), t\rightarrow 0} = \frac{\alpha_s}{4\pi}C_F\times 2\left[\log\left(\bar{z}^2\right) +\gamma_E + 2-2\sqrt{\pi}\bar{z}\right],
\end{eqnarray}
where the last term is consistent with the one-lop result of $\delta m$ we have obtained in eq.~(\ref{eq:deltamresult}).
Notice that the $\overline{\rm MS}$ renormalized result of the Wilson line self energy at one-loop is given by \cite{Izubuchi:2018srq}
\begin{eqnarray}
\mathcal{M}_{(f_1)}^{\overline{\rm MS}} = \frac{\alpha_s}{4\pi}C_F\times 2 \left[2 + 2\gamma_E+ \log\left(\frac{\mu^2z^2}{4}\right)\right].
\end{eqnarray}
Dropping the linearly divergent term from $\mathcal{M}_{(f_1), t\rightarrow 0}$ and comparing with $\mathcal{M}_{(f_1)}^{\overline{\rm MS}} $, gives the contribution from diagram $(f_1)$ for the matching from the gradient-flow scheme to the $\overline{\rm MS}$ scheme in the small flow-time limit,
\begin{eqnarray}
\label{eq:WilsonselfGF2MS}
\mathcal{C}_{(f_1),\, t \rightarrow 0} =  \frac{\alpha_sC_F}{4\pi}\left[-2\log\left(2\mu^2te^{\gamma_E}\right)\right],
\end{eqnarray}
which is consistent with the result $\zeta_{h_v}$ we have obtained eq.~(\ref{eq:Zhvresult}).

%%%%%%%%%%%%%%%%%%%%%%%%%%%%%%%%%%%%%%%%%%%%%%%%%%%%%%%%%%%%%%%%%%%
\subsection{Vertex diagram $(f_2)$}
The amplitude of diagram $(f_2)$ is given by (including its mirror diagram),
\begin{eqnarray}
\mathcal{M}_{(f_2)} =2igC_F \int_0^z d\, s\int \frac{d^d k}{(2\pi)^d} \frac{-i\slashed{k}}{(k^2)^2} \left(ig\gamma\cdot v\right)\, e^{-2t k^2} e^{-i(k\cdot v)s}.
\end{eqnarray}
Since the transverse components of $\slashed{k}$ lead to vanishing momentum integration, we can replace $\slashed{k}$ with $(k\cdot v)(\gamma\cdot v)$ in $\mathcal{M}_{(f_2)}$, therefore, we have
\begin{eqnarray}
\label{eq:resultsdiagramd}
\mathcal{M}_{(f_2)} &=& 2g^2C_F \int \frac{d^d k}{(2\pi)^d} \frac{1}{(k^2)^2} \left(1-e^{-i(k\cdot v) z}\right)\, e^{-2t k^2} \nonumber\\
&=& \frac{\alpha_s}{4\pi}C_F\times 2\left[\log\left(\bar{z}^2\right) + \gamma_E - 1 -\text{Ei}\left(-\bar{z}^2\right) + \frac{1}{\bar{z}^2}\left(1-e^{-\bar{z}^2}\right)\right].
\end{eqnarray}
In the small flow-time limit, we have
\begin{eqnarray}
\mathcal{M}_{(f_2), t\rightarrow 0} &=&\frac{\alpha_s}{4\pi}C_F\times 2\left[\log\left(\frac{z^2}{8t}\right) + \gamma_E - 1 \right].
\end{eqnarray}
Comparing the corresponding result in the $\overline{\rm MS}$ scheme \cite{Izubuchi:2018srq},
\begin{eqnarray}
\mathcal{M}_{(f_2)}^{\overline{\rm MS}} &=&\frac{\alpha_s}{4\pi}C_F\times 2\left[2\gamma_E  +\log\left(\frac{\mu^2z^2}{4}\right) \right],
\end{eqnarray}
we obtain the contribution from diagram $(f_2)$ for the matching from the gradient-flow scheme to the $\overline{\rm MS}$ scheme in the small flow-time limit,
\begin{eqnarray}
\label{eq:quarkvertexGF2MS}
\mathcal{C}_{(f_2), t\rightarrow 0} = \frac{\alpha_s}{4\pi}C_F\left[-2\log\left(2\mu^2te^{\gamma_E}\right)  -2 \right],
\end{eqnarray}
which is consistent with the result of $\zeta_{\psi h_v}$ we have obtained in eq.~(\ref{eq:Zqhvresults}).

\subsection{Diagram $(f_3)$}
The amplitude of diagram $(f_3)$ can be expressed as,
\begin{eqnarray}
\mathcal{M}_{(f_3)} &=&(ig)^2\tilde{\mu}^{2\epsilon} C_F \int \frac{d^d k}{(2\pi)^d} \frac{\delta^{\mu\nu}\gamma^\mu(-i\slashed{k})\gamma^\alpha(-i\slashed{k})\gamma^\nu}{(k^2)^3}e^{-2t k^2}e^{-i(k\cdot v)z}\nonumber\\
&=& g^2C_F \left(d-2\right)\tilde{\mu}^{2\epsilon} \int \frac{d^d k}{(2\pi)^d} \frac{k^2\gamma^\alpha -2k^\alpha\slashed{k}}{(k^2)^3}e^{-2t k^2}e^{-i(k\cdot v)z}.
\end{eqnarray}
If $\alpha$ is parallel to $v$, dropping the tree-level factor $v^{\alpha}\gamma\cdot v $, we have
\begin{eqnarray}
\label{eq:resultsdiagrama0}
\mathcal{M}_{(f_3),\, \|}
&=& g^2C_F \left(d-2\right)\tilde{\mu}^{2\epsilon}  \int \frac{d^d k}{(2\pi)^d} \frac{\left[k^2-2(k\cdot v)^2\right]}{(k^2)^3}e^{-2t k^2}e^{-i(k\cdot v)z}\nonumber\\
&=&\frac{\alpha_s}{4\pi}C_F\bigg[-\frac{1}{\epsilon_{\rm IR}} - \log\left(\frac{\mu^2z^2}{4}\right) -2\gamma_E +3 \nonumber\\
&&\, \, \, \, \, \, \, \,\, \, \, \, \, \, \, \, \,  \, \, \,  +\text{Ei}(-\bar{z}^2) +\frac{3}{\bar{z}^4}\left(1-e^{-\bar{z}^2}\right) - \frac{1}{\bar{z}^2}\left(4-e^{-\bar{z}^2}\right)\bigg].
\end{eqnarray}
If $\alpha$ is transverse to $v$,  dropping the tree-level factor $\gamma^{\alpha}_{\perp}$, we have
\begin{eqnarray}
\label{eq:resultsdiagramaalpha}
\mathcal{M}_{(f_3),\, \perp}&=& g^2C_F \frac{d-2}{d-1}\tilde{\mu}^{2\epsilon}  \int \frac{d^d k}{(2\pi)^d} \frac{\left(d-3\right)k^2+2(k\cdot v)^2}{(k^2)^3}e^{-2t k^2}e^{-i(k\cdot v)z}\nonumber\\
&=&\frac{\alpha_s}{4\pi}C_F\bigg[-\frac{1}{\epsilon_{\rm IR}} - \log\left(\frac{\mu^2z^2}{4}\right) -2\gamma_E +1 \nonumber\\
&&\, \, \, \, \, \, \, \,\, \, \, \, \, \, \, \, \,  \, \, \,  +\text{Ei}(-\bar{z}^2) - \frac{1}{\bar{z}^4}\left(1-e^{-\bar{z}^2}\right) + \frac{1}{\bar{z}^2}
e^{-\bar{z}^2}\bigg].
\end{eqnarray}
In the small flow-time limit, we have
\begin{eqnarray}
\mathcal{M}_{(f_3),\, \|}^{t\rightarrow 0} &= &\frac{\alpha_s}{4\pi}C_F\bigg[-\frac{1}{\epsilon_{\rm IR}} - \log\left(\frac{\mu^2z^2}{4}\right) -2\gamma_E +3 \bigg],\\
\mathcal{M}_{(f_3),\, \perp}^{t\rightarrow 0} &= &\frac{\alpha_s}{4\pi}C_F\bigg[-\frac{1}{\epsilon_{\rm IR}} - \log\left(\frac{\mu^2z^2}{4}\right) -2\gamma_E +1 \bigg].
\end{eqnarray}
Noticing the corresponding $\overline{\rm MS}$ results are given by \cite{Izubuchi:2018srq},
\begin{eqnarray}
\mathcal{M}_{(f_3), \|}^{\overline{\rm MS}} &= &\frac{\alpha_s}{4\pi}C_F\bigg[-\frac{1}{\epsilon_{\rm IR}} - \log\left(\frac{\mu^2z^2}{4}\right) -2\gamma_E +3 \bigg],\\
\mathcal{M}_{(f_3), \perp}^{\overline{\rm MS}} &= &\frac{\alpha_s}{4\pi}C_F\bigg[-\frac{1}{\epsilon_{\rm IR}} - \log\left(\frac{\mu^2z^2}{4}\right) -2\gamma_E +1 \bigg],
\end{eqnarray}
which are exactly the same with the gradient flow results in the small flow-time limit. This indicates that diagram $(f_3)$ does not contribute to the matching from the gradient-flow scheme to 
the $\overline{\rm MS}$ scheme in the small flow-time limit, which supports our argument that we only need the matching for the local current operators stemming from the Wilson-line operators based on the one-dimensional auxiliary-field formalism.

\subsection{Final results}
Adding up the full result of each diagram in figure.~\ref{fig:sqpdf1} and the result of quark self-energy diagram in figure.~\ref{fig:quarkself}, which are given by eq.~(\ref{eq:Wilsonselffull}), eq.~(\ref{eq:resultsdiagramd}), 
eq.~(\ref{eq:resultsdiagrama0}), eq.~(\ref{eq:resultsdiagramaalpha}) and eq.~(\ref{eq:appquarkself}), dropping the linearly divergent term through ``mass” renormalization, we obtain the renormalized full result in the gradient-flow scheme
\begin{eqnarray}
\label{eq:quarkquasiGFfull}
\mathcal{M}_{q} = \frac{\alpha_s}{4\pi}C_F\left[a_{\Gamma}  + 3\log\left(\bar{z}^2e^{\gamma_E}\right)-\log(432)-4e^{-\bar{z}^2} - 3\text{Ei}(-\bar{z}^2) -4\sqrt{\pi}\bar{z}\text{erf}(\bar{z}) + 4\sqrt{\pi}\bar{z} \right],\nonumber\\
\end{eqnarray}
where
\begin{eqnarray}
a_{\gamma\cdot v}&=& 5 +\frac{3}{\bar{z}^4}\left(1-e^{-\bar{z}^2}\right) -\frac{1}{\bar{z}^2}\left(2+e^{-\bar{z}^2}\right),\\
a_{\gamma^{\alpha}_\perp}&=& 3 - \frac{1}{\bar{z}^4}\left(1-e^{-\bar{z}^2}\right) + \frac{1}{\bar{z}^2}\left(2-e^{-\bar{z}^2}\right).
\end{eqnarray}
The $\overline{\rm MS}$ renormalized  total result is 
\begin{eqnarray}
\mathcal{M}_{q}^{\overline{\rm MS}} = \frac{\alpha_s}{4\pi}C_F\left[a^\prime_{\Gamma}+ 3\log\left(\frac{z^2\mu^2}{4}\right)+6\gamma_E\right],
\end{eqnarray}
with $a^\prime_{\gamma\cdot v} = 7,\, a^\prime_{\gamma^{\alpha}_\perp} = 5$.

Then, we have the full matching coefficient, 
\begin{eqnarray}
\mathcal{C}_q (t, \mu, z)& =& 1 +\frac{\alpha_s}{4\pi}C_F\bigg[a_\Gamma  - a^\prime_\Gamma - 3\log\left(2\mu^2te^{\gamma_E}\right)-\log(432) - 4e^{-\bar{z}^2}\nonumber\\
&&\, \, \, \, \, \, \, \, \, \, \, \, \, \, \, \, \, \, \, \, \, \, \, \,\, \, \, \, \, \, - 3 \text{Ei}(-\bar{z}^2) -4\sqrt{\pi}\bar{z}\text{erf}(\bar{z}) + 4\sqrt{\pi}\bar{z}\bigg] +O(\alpha_s^2),
\end{eqnarray}
which gives the small flow-time limit,
\begin{eqnarray}
\mathcal{C}_{q,\, t\rightarrow 0}(t, \mu) = 1 -  \frac{\alpha_s}{4\pi}C_F\left[3\log\left(2\mu^2te^{\gamma_E}\right)+2+\log(432)\right],
\end{eqnarray}
which is indeed independent of $z$ and $\Gamma$.
It is straightforward to check that $\mathcal{C}_{q,\, t\rightarrow 0}(t, \mu) = c_{\psi h_v}^2$ as expected (see eq.~(\ref{eq:cpsihvfinal})).
\begin{figure}
\centering
\includegraphics[width =0.72\textwidth]{images/plot.pdf}
\caption{The ratio of the one-loop correction result with full flow time dependence and the one-loop correction result in the small flow-time limit: ${\rm Ratio} =\frac{\mathcal{C}_q (t, \mu, z)-1}{\mathcal{C}_{q,\, t\rightarrow 0}(t, \mu) -1}$. The renormalization scale $\mu$ is chosen so that $\log\left(2t\mu^2e^{\gamma_E}\right)=0$.}
\label{fig:plot}
\end{figure}

To see how large the finite flow time effect will be, we plot the ratio of the one-loop corrections with full flow time dependence and the one-loop corrections in the small flow-time limit in figure.~\ref{fig:plot}.
As it is shown in figure.~\ref{fig:plot}, the finite flow time effect is very small as long as the flow radius $r_F = \sqrt{8t}$ is smaller than the distance $z$ ($\bar{z}>1$), which is usually satisfied in lattice gradient flow computations.






\vspace{1em}




\begingroup\raggedright
\begin{thebibliography}{10}
\bibitem{Polyakov:1980ca}
A.~M. Polyakov, \emph{{Gauge Fields as Rings of Glue}},
  \href{https://doi.org/10.1016/0550-3213(80)90507-6}{\emph{Nucl. Phys. B}
  {\bfseries 164} (1980) 171}.

\bibitem{Gervais:1979fv}
J.-L. Gervais and A.~Neveu, \emph{{The Slope of the Leading Regge Trajectory in
  Quantum Chromodynamics}},
  \href{https://doi.org/10.1016/0550-3213(80)90397-1}{\emph{Nucl. Phys. B}
  {\bfseries 163} (1980) 189}.

\bibitem{Dotsenko:1979wb}
V.~S. Dotsenko and S.~N. Vergeles, \emph{{Renormalizability of Phase Factors in
  the Nonabelian Gauge Theory}},
  \href{https://doi.org/10.1016/0550-3213(80)90103-0}{\emph{Nucl. Phys. B}
  {\bfseries 169} (1980) 527}.

\bibitem{Craigie:1980qs}
N.~S. Craigie and H.~Dorn, \emph{{On the Renormalization and Short Distance
  Properties of Hadronic Operators in {QCD}}},
  \href{https://doi.org/10.1016/0550-3213(81)90372-2}{\emph{Nucl. Phys. B}
  {\bfseries 185} (1981) 204}.

\bibitem{Dorn:1986dt}
H.~Dorn, \emph{{Renormalization of Path Ordered Phase Factors and Related
  Hadron Operators in Gauge Field Theories}},
  \href{https://doi.org/10.1002/prop.19860340104}{\emph{Fortsch. Phys.}
  {\bfseries 34} (1986) 11}.

\bibitem{Arefeva:1980zd}
I.~Y. Arefeva, \emph{{QUANTUM CONTOUR FIELD EQUATIONS}},
  \href{https://doi.org/10.1016/0370-2693(80)90529-8}{\emph{Phys. Lett. B}
  {\bfseries 93} (1980) 347}.

\bibitem{Arefeva:1980uy}
I.~Y. Arefeva, \emph{{Elimination of Divergences in an Integral Formulation of
  {Yang-Mills} Theory}}, {\emph{JETP Lett.} {\bfseries 31} (1980) 393}.

\bibitem{Samuel:1978iy}
S.~Samuel, \emph{{COLOR ZITTERBEWEGUNG}},
  \href{https://doi.org/10.1016/0550-3213(79)90005-1}{\emph{Nucl. Phys. B}
  {\bfseries 149} (1979) 517}.

\bibitem{Brandt:1978wc}
R.~A. Brandt, F.~Neri and D.~Zwanziger, \emph{{Lorentz Invariance From
  Classical Particle Paths in Quantum Field Theory of Electric and Magnetic
  Charge}}, \href{https://doi.org/10.1103/PhysRevD.19.1153}{\emph{Phys. Rev. D}
  {\bfseries 19} (1979) 1153}.

\bibitem{Ji:2013dva}
X.~Ji, \emph{{Parton Physics on a Euclidean Lattice}},
  \href{https://doi.org/10.1103/PhysRevLett.110.262002}{\emph{Phys. Rev. Lett.}
  {\bfseries 110} (2013) 262002}
  [\href{https://arxiv.org/abs/1305.1539}{{\ttfamily 1305.1539}}].

\bibitem{Radyushkin:2017cyf}
A.~V. Radyushkin, \emph{{Quasi-parton distribution functions, momentum
  distributions, and pseudo-parton distribution functions}},
  \href{https://doi.org/10.1103/PhysRevD.96.034025}{\emph{Phys. Rev. D}
  {\bfseries 96} (2017) 034025}
  [\href{https://arxiv.org/abs/1705.01488}{{\ttfamily 1705.01488}}].

\bibitem{Pineda:1997bj}
A.~Pineda and J.~Soto, \emph{{Effective field theory for ultrasoft momenta in
  NRQCD and NRQED}},
  \href{https://doi.org/10.1016/S0920-5632(97)01102-X}{\emph{Nucl. Phys. B
  Proc. Suppl.} {\bfseries 64} (1998) 428}
  [\href{https://arxiv.org/abs/hep-ph/9707481}{{\ttfamily hep-ph/9707481}}].

\bibitem{Brambilla:1999xf}
N.~Brambilla, A.~Pineda, J.~Soto and A.~Vairo, \emph{{Potential NRQCD: An
  Effective theory for heavy quarkonium}},
  \href{https://doi.org/10.1016/S0550-3213(99)00693-8}{\emph{Nucl. Phys. B}
  {\bfseries 566} (2000) 275}
  [\href{https://arxiv.org/abs/hep-ph/9907240}{{\ttfamily hep-ph/9907240}}].

\bibitem{Brambilla:2004jw}
N.~Brambilla, A.~Pineda, J.~Soto and A.~Vairo, \emph{{Effective Field Theories
  for Heavy Quarkonium}},
  \href{https://doi.org/10.1103/RevModPhys.77.1423}{\emph{Rev. Mod. Phys.}
  {\bfseries 77} (2005) 1423}
  [\href{https://arxiv.org/abs/hep-ph/0410047}{{\ttfamily hep-ph/0410047}}].

\bibitem{Brambilla:2001xy}
N.~Brambilla, D.~Eiras, A.~Pineda, J.~Soto and A.~Vairo, \emph{{New predictions
  for inclusive heavy quarkonium P wave decays}},
  \href{https://doi.org/10.1103/PhysRevLett.88.012003}{\emph{Phys. Rev. Lett.}
  {\bfseries 88} (2002) 012003}
  [\href{https://arxiv.org/abs/hep-ph/0109130}{{\ttfamily hep-ph/0109130}}].

\bibitem{Brambilla:2002nu}
N.~Brambilla, D.~Eiras, A.~Pineda, J.~Soto and A.~Vairo, \emph{{Inclusive
  decays of heavy quarkonium to light particles}},
  \href{https://doi.org/10.1103/PhysRevD.67.034018}{\emph{Phys. Rev. D}
  {\bfseries 67} (2003) 034018}
  [\href{https://arxiv.org/abs/hep-ph/0208019}{{\ttfamily hep-ph/0208019}}].

\bibitem{Brambilla:2020ojz}
N.~Brambilla, H.~S. Chung and A.~Vairo, \emph{{Inclusive Hadroproduction of
  $P$-Wave Heavy Quarkonia in Potential Nonrelativistic QCD}},
  \href{https://doi.org/10.1103/PhysRevLett.126.082003}{\emph{Phys. Rev. Lett.}
  {\bfseries 126} (2021) 082003}
  [\href{https://arxiv.org/abs/2007.07613}{{\ttfamily 2007.07613}}].

\bibitem{Brambilla:2021abf}
N.~Brambilla, H.~S. Chung and A.~Vairo, \emph{{Inclusive production of heavy
  quarkonia in pNRQCD}},
  \href{https://doi.org/10.1007/JHEP09(2021)032}{\emph{JHEP} {\bfseries 09}
  (2021) 032} [\href{https://arxiv.org/abs/2106.09417}{{\ttfamily
  2106.09417}}].

\bibitem{Brambilla:2022rjd}
N.~Brambilla, H.~S. Chung, A.~Vairo and X.-P. Wang, \emph{{Production and
  polarization of S-wave quarkonia in potential nonrelativistic QCD}},
  \href{https://doi.org/10.1103/PhysRevD.105.L111503}{\emph{Phys. Rev. D}
  {\bfseries 105} (2022) L111503}
  [\href{https://arxiv.org/abs/2203.07778}{{\ttfamily 2203.07778}}].

\bibitem{Brambilla:2022ayc}
N.~Brambilla, H.~S. Chung, A.~Vairo and X.-P. Wang, \emph{{Inclusive production
  of J/\ensuremath{\psi}, \ensuremath{\psi}(2S), and \ensuremath{\Upsilon}
  states in pNRQCD}},
  \href{https://doi.org/10.1007/JHEP03(2023)242}{\emph{JHEP} {\bfseries 03}
  (2023) 242} [\href{https://arxiv.org/abs/2210.17345}{{\ttfamily
  2210.17345}}].

\bibitem{Braun:2007wv}
V.~Braun and D.~M\"uller, \emph{{Exclusive processes in position space and the
  pion distribution amplitude}},
  \href{https://doi.org/10.1140/epjc/s10052-008-0608-4}{\emph{Eur. Phys. J. C}
  {\bfseries 55} (2008) 349} [\href{https://arxiv.org/abs/0709.1348}{{\ttfamily
  0709.1348}}].

\bibitem{Ji:2014gla}
X.~Ji, \emph{{Parton Physics from Large-Momentum Effective Field Theory}},
  \href{https://doi.org/10.1007/s11433-014-5492-3}{\emph{Sci. China Phys. Mech.
  Astron.} {\bfseries 57} (2014) 1407}
  [\href{https://arxiv.org/abs/1404.6680}{{\ttfamily 1404.6680}}].

\bibitem{Izubuchi:2018srq}
T.~Izubuchi, X.~Ji, L.~Jin, I.~W. Stewart and Y.~Zhao, \emph{{Factorization
  Theorem Relating Euclidean and Light-Cone Parton Distributions}},
  \href{https://doi.org/10.1103/PhysRevD.98.056004}{\emph{Phys. Rev. D}
  {\bfseries 98} (2018) 056004}
  [\href{https://arxiv.org/abs/1801.03917}{{\ttfamily 1801.03917}}].

\bibitem{Ji:2020ect}
X.~Ji, Y.-S. Liu, Y.~Liu, J.-H. Zhang and Y.~Zhao, \emph{{Large-momentum
  effective theory}},
  \href{https://doi.org/10.1103/RevModPhys.93.035005}{\emph{Rev. Mod. Phys.}
  {\bfseries 93} (2021) 035005}
  [\href{https://arxiv.org/abs/2004.03543}{{\ttfamily 2004.03543}}].

\bibitem{DiGiacomo:2000irz}
A.~Di~Giacomo, H.~G. Dosch, V.~I. Shevchenko and Y.~A. Simonov, \emph{{Field
  correlators in QCD: Theory and applications}},
  \href{https://doi.org/10.1016/S0370-1573(02)00140-0}{\emph{Phys. Rept.}
  {\bfseries 372} (2002) 319}
  [\href{https://arxiv.org/abs/hep-ph/0007223}{{\ttfamily hep-ph/0007223}}].

\bibitem{Narayanan:2006rf}
R.~Narayanan and H.~Neuberger, \emph{{Infinite N phase transitions in continuum
  Wilson loop operators}},
  \href{https://doi.org/10.1088/1126-6708/2006/03/064}{\emph{JHEP} {\bfseries
  03} (2006) 064} [\href{https://arxiv.org/abs/hep-th/0601210}{{\ttfamily
  hep-th/0601210}}].

\bibitem{Luscher:2009eq}
M.~L{\"u}scher, \emph{{Trivializing maps, the Wilson flow and the HMC
  algorithm}}, \href{https://doi.org/10.1007/s00220-009-0953-7}{\emph{Commun.
  Math. Phys.} {\bfseries 293} (2010) 899}
  [\href{https://arxiv.org/abs/0907.5491}{{\ttfamily 0907.5491}}].

\bibitem{Luscher:2010iy}
M.~L\"uscher, \emph{{Properties and uses of the Wilson flow in lattice QCD}},
  \href{https://doi.org/10.1007/JHEP08(2010)071}{\emph{JHEP} {\bfseries 08}
  (2010) 071} [\href{https://arxiv.org/abs/1006.4518}{{\ttfamily 1006.4518}}].

\bibitem{Luscher:2011bx}
M.~Luscher and P.~Weisz, \emph{{Perturbative analysis of the gradient flow in
  non-abelian gauge theories}},
  \href{https://doi.org/10.1007/JHEP02(2011)051}{\emph{JHEP} {\bfseries 02}
  (2011) 051} [\href{https://arxiv.org/abs/1101.0963}{{\ttfamily 1101.0963}}].

\bibitem{Luscher:2013cpa}
M.~Luscher, \emph{{Chiral symmetry and the Yang--Mills gradient flow}},
  \href{https://doi.org/10.1007/JHEP04(2013)123}{\emph{JHEP} {\bfseries 04}
  (2013) 123} [\href{https://arxiv.org/abs/1302.5246}{{\ttfamily 1302.5246}}].

\bibitem{Luscher:2013vga}
M.~L\"uscher, \emph{{Future applications of the Yang\textendash{}Mills gradient
  flow in lattice QCD}}, \href{https://doi.org/10.22323/1.187.0016}{\emph{PoS}
  {\bfseries LATTICE2013} (2014) 016}
  [\href{https://arxiv.org/abs/1308.5598}{{\ttfamily 1308.5598}}].

\bibitem{Borsanyi:2012zs}
S.~Borsanyi et~al., \emph{{High-precision scale setting in lattice QCD}},
  \href{https://doi.org/10.1007/JHEP09(2012)010}{\emph{JHEP} {\bfseries 09}
  (2012) 010} [\href{https://arxiv.org/abs/1203.4469}{{\ttfamily 1203.4469}}].

\bibitem{Sommer:2014mea}
R.~Sommer, \emph{{Scale setting in lattice QCD}},
  \href{https://doi.org/10.22323/1.187.0015}{\emph{PoS} {\bfseries LATTICE2013}
  (2014) 015} [\href{https://arxiv.org/abs/1401.3270}{{\ttfamily 1401.3270}}].

\bibitem{Suzuki:2013gza}
H.~Suzuki, \emph{{Energy\textendash{}momentum tensor from the
  Yang\textendash{}Mills gradient flow}},
  \href{https://doi.org/10.1093/ptep/ptt059}{\emph{PTEP} {\bfseries 2013}
  (2013) 083B03} [\href{https://arxiv.org/abs/1304.0533}{{\ttfamily
  1304.0533}}].

\bibitem{Makino:2014taa}
H.~Makino and H.~Suzuki, \emph{{Lattice energy\textendash{}momentum tensor from
  the Yang\textendash{}Mills gradient flow\textemdash{}inclusion of fermion
  fields}}, \href{https://doi.org/10.1093/ptep/ptu070}{\emph{PTEP} {\bfseries
  2014} (2014) 063B02} [\href{https://arxiv.org/abs/1403.4772}{{\ttfamily
  1403.4772}}].

\bibitem{Harlander:2018zpi}
R.~V. Harlander, Y.~Kluth and F.~Lange, \emph{{The two-loop
  energy\textendash{}momentum tensor within the gradient-flow formalism}},
  \href{https://doi.org/10.1140/epjc/s10052-018-6415-7}{\emph{Eur. Phys. J. C}
  {\bfseries 78} (2018) 944}
  [\href{https://arxiv.org/abs/1808.09837}{{\ttfamily 1808.09837}}].

\bibitem{FlavourLatticeAveragingGroupFLAG:2021npn}
{\scshape Flavour Lattice Averaging Group (FLAG)} collaboration, \emph{{FLAG
  Review 2021}},
  \href{https://doi.org/10.1140/epjc/s10052-022-10536-1}{\emph{Eur. Phys. J. C}
  {\bfseries 82} (2022) 869}
  [\href{https://arxiv.org/abs/2111.09849}{{\ttfamily 2111.09849}}].

\bibitem{Leino:2021vop}
V.~Leino, N.~Brambilla, J.~Mayer-Steudte and A.~Vairo, \emph{{The static force
  from generalized Wilson loops using gradient flow}},
  \href{https://doi.org/10.1051/epjconf/202225804009}{\emph{EPJ Web Conf.}
  {\bfseries 258} (2022) 04009}
  [\href{https://arxiv.org/abs/2111.10212}{{\ttfamily 2111.10212}}].

\bibitem{Mayer-Steudte:2022uih}
J.~Mayer-Steudte, N.~Brambilla, V.~Leino and A.~Vairo, \emph{{Implications of
  gradient flow on the static force}},
  \href{https://doi.org/10.22323/1.430.0353}{\emph{PoS} {\bfseries LATTICE2022}
  (2023) 353} [\href{https://arxiv.org/abs/2212.12400}{{\ttfamily
  2212.12400}}].

\bibitem{Leino:2022kgj}
V.~Leino, N.~Brambilla, J.~Mayer-Steudte and P.~Petreczky, \emph{{Heavy quark
  diffusion coefficient with gradient flow}},
  \href{https://doi.org/10.22323/1.430.0183}{\emph{PoS} {\bfseries LATTICE2022}
  (2023) 183} [\href{https://arxiv.org/abs/2212.10941}{{\ttfamily
  2212.10941}}].

\bibitem{Monahan:2016bvm}
C.~Monahan and K.~Orginos, \emph{{Quasi parton distributions and the gradient
  flow}}, \href{https://doi.org/10.1007/JHEP03(2017)116}{\emph{JHEP} {\bfseries
  03} (2017) 116} [\href{https://arxiv.org/abs/1612.01584}{{\ttfamily
  1612.01584}}].

\bibitem{Monahan:2017hpu}
C.~Monahan, \emph{{Smeared quasidistributions in perturbation theory}},
  \href{https://doi.org/10.1103/PhysRevD.97.054507}{\emph{Phys. Rev. D}
  {\bfseries 97} (2018) 054507}
  [\href{https://arxiv.org/abs/1710.04607}{{\ttfamily 1710.04607}}].

\bibitem{Artz:2019bpr}
J.~Artz, R.~V. Harlander, F.~Lange, T.~Neumann and M.~Prausa, \emph{{Results
  and techniques for higher order calculations within the gradient-flow
  formalism}}, \href{https://doi.org/10.1007/JHEP06(2019)121}{\emph{JHEP}
  {\bfseries 06} (2019) 121}
  [\href{https://arxiv.org/abs/1905.00882}{{\ttfamily 1905.00882}}].

\bibitem{Rizik:2020naq}
{\scshape SymLat} collaboration, \emph{{Short flow-time coefficients of
  $CP$-violating operators}},
  \href{https://doi.org/10.1103/PhysRevD.102.034509}{\emph{Phys. Rev. D}
  {\bfseries 102} (2020) 034509}
  [\href{https://arxiv.org/abs/2005.04199}{{\ttfamily 2005.04199}}].

\bibitem{Suzuki:2020zue}
A.~Suzuki, Y.~Taniguchi, H.~Suzuki and K.~Kanaya, \emph{{Four quark operators
  for kaon bag parameter with gradient flow}},
  \href{https://doi.org/10.1103/PhysRevD.102.034508}{\emph{Phys. Rev. D}
  {\bfseries 102} (2020) 034508}
  [\href{https://arxiv.org/abs/2006.06999}{{\ttfamily 2006.06999}}].

\bibitem{Harlander:2020duo}
R.~V. Harlander, F.~Lange and T.~Neumann, \emph{{Hadronic vacuum polarization
  using gradient flow}},
  \href{https://doi.org/10.1007/JHEP08(2020)109}{\emph{JHEP} {\bfseries 08}
  (2020) 109} [\href{https://arxiv.org/abs/2007.01057}{{\ttfamily
  2007.01057}}].

\bibitem{Boers:2020uqr}
M.~Boers and E.~Pallante, \emph{{Conserved vector current in QCD-like theories
  and the gradient flow}},
  \href{https://doi.org/10.1007/JHEP10(2020)034}{\emph{JHEP} {\bfseries 10}
  (2020) 034} [\href{https://arxiv.org/abs/2007.02121}{{\ttfamily
  2007.02121}}].

\bibitem{Shindler:2021bcx}
A.~Shindler, \emph{{Flavor-diagonal CP violation: the electric dipole moment}},
  \href{https://doi.org/10.1140/epja/s10050-021-00421-y}{\emph{Eur. Phys. J. A}
  {\bfseries 57} (2021) 128}.

\bibitem{Harlander:2022tgk}
R.~V. Harlander and F.~Lange, \emph{{Effective electroweak Hamiltonian in the
  gradient-flow formalism}},
  \href{https://doi.org/10.1103/PhysRevD.105.L071504}{\emph{Phys. Rev. D}
  {\bfseries 105} (2022) L071504}
  [\href{https://arxiv.org/abs/2201.08618}{{\ttfamily 2201.08618}}].

\bibitem{Mereghetti:2021nkt}
E.~Mereghetti, C.~J. Monahan, M.~D. Rizik, A.~Shindler and P.~Stoffer,
  \emph{{One-loop matching for quark dipole operators in a gradient-flow
  scheme}}, \href{https://doi.org/10.1007/JHEP04(2022)050}{\emph{JHEP}
  {\bfseries 04} (2022) 050}
  [\href{https://arxiv.org/abs/2111.11449}{{\ttfamily 2111.11449}}].

\bibitem{Harlander:2022vgf}
R.~Harlander, M.~D. Rizik, J.~Borgulat and A.~Shindler, \emph{{Two-loop
  matching of the chromo-magnetic dipole operator with the gradient flow}},
  \href{https://doi.org/10.22323/1.430.0313}{\emph{PoS} {\bfseries LATTICE2022}
  (2023) 313} [\href{https://arxiv.org/abs/2212.09824}{{\ttfamily
  2212.09824}}].

\bibitem{Brambilla:2021egm}
N.~Brambilla, H.~S. Chung, A.~Vairo and X.-P. Wang, \emph{{QCD static force in
  gradient flow}}, \href{https://doi.org/10.1007/JHEP01(2022)184}{\emph{JHEP}
  {\bfseries 01} (2022) 184}
  [\href{https://arxiv.org/abs/2111.07811}{{\ttfamily 2111.07811}}].

\bibitem{Crosas:2023anw}
O.~L. Crosas, C.~J. Monahan, M.~D. Rizik, A.~Shindler and P.~Stoffer,
  \emph{{One-loop matching of the $CP$-odd three-gluon operator to the gradient
  flow}},  \href{https://arxiv.org/abs/2308.16221}{{\ttfamily 2308.16221}}.

\bibitem{Eichten:1989zv}
E.~Eichten and B.~R. Hill, \emph{{An Effective Field Theory for the Calculation
  of Matrix Elements Involving Heavy Quarks}},
  \href{https://doi.org/10.1016/0370-2693(90)92049-O}{\emph{Phys. Lett. B}
  {\bfseries 234} (1990) 511}.

\bibitem{Bigi:1994em}
I.~I.~Y. Bigi, M.~A. Shifman, N.~G. Uraltsev and A.~I. Vainshtein, \emph{{The
  Pole mass of the heavy quark. Perturbation theory and beyond}},
  \href{https://doi.org/10.1103/PhysRevD.50.2234}{\emph{Phys. Rev. D}
  {\bfseries 50} (1994) 2234}
  [\href{https://arxiv.org/abs/hep-ph/9402360}{{\ttfamily hep-ph/9402360}}].

\bibitem{Beneke:1994sw}
M.~Beneke and V.~M. Braun, \emph{{Heavy quark effective theory beyond
  perturbation theory: Renormalons, the pole mass and the residual mass term}},
  \href{https://doi.org/10.1016/0550-3213(94)90314-X}{\emph{Nucl. Phys. B}
  {\bfseries 426} (1994) 301}
  [\href{https://arxiv.org/abs/hep-ph/9402364}{{\ttfamily hep-ph/9402364}}].

\bibitem{Braun:2020ymy}
V.~M. Braun, K.~G. Chetyrkin and B.~A. Kniehl, \emph{{Renormalization of parton
  quasi-distributions beyond the leading order: spacelike vs. timelike}},
  \href{https://doi.org/10.1007/JHEP07(2020)161}{\emph{JHEP} {\bfseries 07}
  (2020) 161} [\href{https://arxiv.org/abs/2004.01043}{{\ttfamily
  2004.01043}}].

\bibitem{Dorn:1981wa}
H.~Dorn, D.~Robaschik and E.~Wieczorek, \emph{{RENORMALIZATION AND SHORT
  DISTANCE PROPERTIES OF GAUGE INVARIANT GLUONIUM AND HADRON OPERATORS}},
  \href{https://doi.org/10.1002/andp.19834950208}{\emph{Annalen Phys.}
  {\bfseries 40} (1983) 166}.

\bibitem{Zhang:2018diq}
J.-H. Zhang, X.~Ji, A.~Sch\"afer, W.~Wang and S.~Zhao, \emph{{Accessing Gluon
  Parton Distributions in Large Momentum Effective Theory}},
  \href{https://doi.org/10.1103/PhysRevLett.122.142001}{\emph{Phys. Rev. Lett.}
  {\bfseries 122} (2019) 142001}
  [\href{https://arxiv.org/abs/1808.10824}{{\ttfamily 1808.10824}}].

\bibitem{Wang:2019tgg}
W.~Wang, J.-H. Zhang, S.~Zhao and R.~Zhu, \emph{{Complete matching for
  quasidistribution functions in large momentum effective theory}},
  \href{https://doi.org/10.1103/PhysRevD.100.074509}{\emph{Phys. Rev. D}
  {\bfseries 100} (2019) 074509}
  [\href{https://arxiv.org/abs/1904.00978}{{\ttfamily 1904.00978}}].

\bibitem{Green:2017xeu}
J.~Green, K.~Jansen and F.~Steffens, \emph{{Nonperturbative Renormalization of
  Nonlocal Quark Bilinears for Parton Quasidistribution Functions on the
  Lattice Using an Auxiliary Field}},
  \href{https://doi.org/10.1103/PhysRevLett.121.022004}{\emph{Phys. Rev. Lett.}
  {\bfseries 121} (2018) 022004}
  [\href{https://arxiv.org/abs/1707.07152}{{\ttfamily 1707.07152}}].

\bibitem{Ji:2017oey}
X.~Ji, J.-H. Zhang and Y.~Zhao, \emph{{Renormalization in Large Momentum
  Effective Theory of Parton Physics}},
  \href{https://doi.org/10.1103/PhysRevLett.120.112001}{\emph{Phys. Rev. Lett.}
  {\bfseries 120} (2018) 112001}
  [\href{https://arxiv.org/abs/1706.08962}{{\ttfamily 1706.08962}}].

\bibitem{Ishikawa:2017faj}
T.~Ishikawa, Y.-Q. Ma, J.-W. Qiu and S.~Yoshida, \emph{{Renormalizability of
  quasiparton distribution functions}},
  \href{https://doi.org/10.1103/PhysRevD.96.094019}{\emph{Phys. Rev. D}
  {\bfseries 96} (2017) 094019}
  [\href{https://arxiv.org/abs/1707.03107}{{\ttfamily 1707.03107}}].

\bibitem{Li:2018tpe}
Z.-Y. Li, Y.-Q. Ma and J.-W. Qiu, \emph{{Multiplicative Renormalizability of
  Operators defining Quasiparton Distributions}},
  \href{https://doi.org/10.1103/PhysRevLett.122.062002}{\emph{Phys. Rev. Lett.}
  {\bfseries 122} (2019) 062002}
  [\href{https://arxiv.org/abs/1809.01836}{{\ttfamily 1809.01836}}].

\bibitem{Stefanis:1983ke}
N.~G. Stefanis, \emph{{GAUGE INVARIANT QUARK TWO POINT GREEN'S FUNCTION THROUGH
  CONNECTOR INSERTION TO O (alpha-s)}},
  \href{https://doi.org/10.1007/BF02902597}{\emph{Nuovo Cim. A} {\bfseries 83}
  (1984) 205}.

\bibitem{Stefanis:2012if}
N.~G. Stefanis, \emph{{Worldline techniques and QCD observables}},
  \href{https://doi.org/10.5506/APhysPolBSupp.6.71}{\emph{Acta Phys. Polon.
  Supp.} {\bfseries 6} (2013) 71}
  [\href{https://arxiv.org/abs/1211.7218}{{\ttfamily 1211.7218}}].

\bibitem{Brambilla:2000gk}
N.~Brambilla, A.~Pineda, J.~Soto and A.~Vairo, \emph{{The QCD potential at
  O(1/m)}}, \href{https://doi.org/10.1103/PhysRevD.63.014023}{\emph{Phys. Rev.
  D} {\bfseries 63} (2001) 014023}
  [\href{https://arxiv.org/abs/hep-ph/0002250}{{\ttfamily hep-ph/0002250}}].

\bibitem{Pineda:2000sz}
A.~Pineda and A.~Vairo, \emph{{The QCD potential at O (1 / $m^{2)}$ : Complete
  spin dependent and spin independent result}},
  \href{https://doi.org/10.1103/PhysRevD.64.039902}{\emph{Phys. Rev. D}
  {\bfseries 63} (2001) 054007}
  [\href{https://arxiv.org/abs/hep-ph/0009145}{{\ttfamily hep-ph/0009145}}].

\bibitem{Eller:2021qpp}
A.~M. Eller, \emph{{The Color-Electric Field Correlator under Gradient Flow at
  next-to-leading Order in Quantum Chromodynamics}}, Ph.D. thesis, Tech. U.,
  Dortmund (main), Darmstadt, Tech. Hochsch., 2021.
\newblock 10.26083/tuprints-00017610.

\bibitem{Dorn:1980hs}
H.~Dorn and E.~Wieczorek, \emph{{Renormalization and Short Distance Properties
  of String Type Equations in {QCD}}},
  \href{https://doi.org/10.1007/BF01554111}{\emph{Z. Phys. C} {\bfseries 9}
  (1981) 49}.

\bibitem{Eichten:1990vp}
E.~Eichten and B.~R. Hill, \emph{{STATIC EFFECTIVE FIELD THEORY: 1/m
  CORRECTIONS}},
  \href{https://doi.org/10.1016/0370-2693(90)91408-4}{\emph{Phys. Lett. B}
  {\bfseries 243} (1990) 427}.

\bibitem{Falk:1990pz}
A.~F. Falk, B.~Grinstein and M.~E. Luke, \emph{{Leading mass corrections to the
  heavy quark effective theory}},
  \href{https://doi.org/10.1016/0550-3213(91)90464-9}{\emph{Nucl. Phys. B}
  {\bfseries 357} (1991) 185}.

\bibitem{Abbott:1981ke}
L.~F. Abbott, \emph{{Introduction to the Background Field Method}}, {\emph{Acta
  Phys. Polon. B} {\bfseries 13} (1982) 33}.

\bibitem{Amoros:1997rx}
G.~Amoros, M.~Beneke and M.~Neubert, \emph{{Two loop anomalous dimension of the
  chromomagnetic moment of a heavy quark}},
  \href{https://doi.org/10.1016/S0370-2693(97)00345-6}{\emph{Phys. Lett. B}
  {\bfseries 401} (1997) 81}
  [\href{https://arxiv.org/abs/hep-ph/9701375}{{\ttfamily hep-ph/9701375}}].

\bibitem{Czarnecki:1997dz}
A.~Czarnecki and A.~G. Grozin, \emph{{HQET chromomagnetic interaction at two
  loops}}, \href{https://doi.org/10.1016/S0370-2693(97)00587-X}{\emph{Phys.
  Lett. B} {\bfseries 405} (1997) 142}
  [\href{https://arxiv.org/abs/hep-ph/9701415}{{\ttfamily hep-ph/9701415}}].

\bibitem{Broadhurst:1991fz}
D.~J. Broadhurst and A.~G. Grozin, \emph{{Two loop renormalization of the
  effective field theory of a static quark}},
  \href{https://doi.org/10.1016/0370-2693(91)90532-U}{\emph{Phys. Lett. B}
  {\bfseries 267} (1991) 105}
  [\href{https://arxiv.org/abs/hep-ph/9908362}{{\ttfamily hep-ph/9908362}}].

\bibitem{Ji:1991pr}
X.-D. Ji and M.~J. Musolf, \emph{{Subleading logarithmic mass dependence in
  heavy meson form-factors}},
  \href{https://doi.org/10.1016/0370-2693(91)91916-J}{\emph{Phys. Lett. B}
  {\bfseries 257} (1991) 409}.

\bibitem{Chetyrkin:2003vi}
K.~G. Chetyrkin and A.~G. Grozin, \emph{{Three loop anomalous dimension of the
  heavy light quark current in HQET}},
  \href{https://doi.org/10.1016/S0550-3213(03)00490-5}{\emph{Nucl. Phys. B}
  {\bfseries 666} (2003) 289}
  [\href{https://arxiv.org/abs/hep-ph/0303113}{{\ttfamily hep-ph/0303113}}].

\bibitem{Constantinou:2017sej}
M.~Constantinou and H.~Panagopoulos, \emph{{Perturbative renormalization of
  quasi-parton distribution functions}},
  \href{https://doi.org/10.1103/PhysRevD.96.054506}{\emph{Phys. Rev. D}
  {\bfseries 96} (2017) 054506}
  [\href{https://arxiv.org/abs/1705.11193}{{\ttfamily 1705.11193}}].

\bibitem{Altenkort:2023eav}
L.~Altenkort, D.~de~la Cruz, O.~Kaczmarek, R.~Larsen, G.~D. Moore, S.~Mukherjee
  et~al., \emph{{Quark Mass Dependence of Heavy Quark Diffusion Coefficient
  from Lattice QCD}},  \href{https://arxiv.org/abs/2311.01525}{{\ttfamily
  2311.01525}}.

\end{thebibliography}\endgroup
